% APPENDIX B: CORE ENGINE CODE SNIPPETS (Massive Expansion)
\chapter{CORE ENGINE CODE SNIPPETS}

This appendix provides the architectural "blueprints" for the Nebula platform. Each snippet includes a technical description of its role in the multi-cloud orchestration engine.

\section{Backend: The Provisioning Controller (FastAPI)}
\begin{lstlisting}[language=Python, caption=Full Provisioning Router]
@router.post("/resources/", response_model=ResourceSchema)
def create_resource(
    request: ResourceCreate, 
    db: Session = Depends(get_db), 
    current_user: User = Depends(get_current_active_user)
):
    # 1. Identity Verification
    if not current_user:
        raise HTTPException(status_code=401, detail="Unauthorized")

    # 2. Persist Intent (Transactional)
    new_resource = Resource(
        user_id=current_user.id,
        name=request.name,
        provider=request.provider,
        type=request.type,
        status="pending",
        configuration=request.variables
    )
    db.add(new_resource)
    db.commit()
    db.refresh(new_resource)

    # 3. Queue Asynchronous Execution
    provision_resource_task.delay(
        resource_id=new_resource.id,
        user_id=current_user.id,
        provider=request.provider,
        module=request.type,
        vars=request.variables
    )
    
    return new_resource
\end{lstlisting}

\section{Security: The Encryption Service}
\begin{lstlisting}[language=Python, caption=AES-128 Fernet Implementation]
from cryptography.fernet import Fernet
import os

def get_vault_engine():
    key = os.getenv("SECRET_KEY")
    if not key:
        raise ValueError("SECRET_KEY not found in environment")
    return Fernet(key)

def encrypt_credential(plain_text: str) -> str:
    f = get_vault_engine()
    return f.encrypt(plain_text.encode()).decode()

def decrypt_credential(cipher_text: str) -> str:
    f = get_vault_engine()
    return f.decrypt(cipher_text.encode()).decode()
\end{lstlisting}

\newpage
% APPENDIX C: CONFIGURATION AND IAC SAMPLES
\chapter{CONFIGURATION AND IAC SAMPLES}

\section{System Orchestration: Docker-Compose}
This file defines the entire multi-container ecosystem of Nebula.

\begin{lstlisting}[language=XML, caption=docker-compose.yml with Worker Cluster]
version: '3.8'

services:
  # 1. Peripheral Services
  db:
    image: postgres:15-alpine
    environment:
      POSTGRES_DB: nebula_db
    volumes:
      - postgres_data:/var/lib/postgresql/data

  redis:
    image: redis:7-alpine
    ports:
      - "6379:6379"

  # 2. Application Core
  api:
    build: ./backend
    command: uvicorn app.main:app --host 0.0.0.0
    environment:
      - DATABASE_URL=postgresql://user:pass@db:5432/nebula_db
      - SECRET_KEY=${SECRET_KEY}
    depends_on:
      - db
      - redis

  worker:
    build: ./backend
    command: celery -A app.tasks worker --loglevel=info
    environment:
      - DATABASE_URL=postgresql://user:pass@db:5432/nebula_db
    depends_on:
      - redis
\end{lstlisting}

\section{IaC Layer: AWS EC2 Terraform Sample}
\begin{lstlisting}[language=XML, caption=AWS VM Module (main.tf)]
resource "aws_instance" "nebula_vm" {
  ami           = var.ami_id
  instance_type = var.instance_type
  
  tags = {
    Name       = var.resource_name
    ManagedBy  = "Nebula-Orchestrator"
  }
}

variable "ami_id" { type = string }
variable "instance_type" { default = "t3.micro" }
variable "resource_name" { type = string }
\end{lstlisting}

\newpage
